% !TEX root = ../main.tex
\clearpage
\chapter{부록}
\label{ch:appendix}

% Demote headings one level under Chapter 8 (sections become 8.1.x).
\renewcommand{\dissertationSubheading}[2][]{%
  \phantomsection
  \subsection{#2}%
  \if\relax\detokenize{#1}\relax\else\label{#1}\fi
}
\renewcommand{\dissertationSubsubheading}[2][]{%
  \phantomsection
  \subsubsection{#2}%
  \if\relax\detokenize{#1}\relax\else\label{#1}\fi
}

\section{평가 파이프라인과 재현성}
\label{ch:appendix-eval}

본 부록은 제~\ref{ch:results}장의 수치를 산출한 평가 파이프라인을 문서화한다. 데이터 출처, 처리 단계, 기본 매개변수, 프록시 공식, 용어집을 포함한다. 여기에 인용된 모든 양은 본문의 표와 그림에도 나타나며, 주장을 해석하기 위해 외부 파일은 필요하지 않다.

\dissertationSubheading[sec:pipeline-stages]{파이프라인 단계}

평가는 Arbitrum One에서 관측 구간 2025-12-01부터 2026-05-31까지(UTC, 양끝 포함) 세 단계로 진행된다.

\begin{enumerate}
\item \textbf{거래 성공 라벨:} 공개 BigQuery 로그 테이블에서 GMX V2 \texttt{PositionDecrease} 이벤트를 디코딩하고, 지갑 신원, 실현 손익(PnL), 포지션 규모 변화, 청산(liquidation) 플래그를 복원한 뒤, 최소 세 번의 close를 가진 지갑을 매칭 코호트로 유지한다($n=5{,}521$; 디코딩된 close $365{,}488$건).
\item \textbf{평판 그래프와 점수:} 코호트 지갑에 대해 ERC-20 \texttt{Approval} 및 \texttt{Transfer} 로그를 월별 배치로 추출한다(쿼리당 dry-run 및 $150$\,GiB 중단 상한). 최신 allowance 보증(endorsement) 그래프와 시간 감쇠 전송 그래프를 구축한 뒤, 거듭제곱 반복(power iteration)으로 EndorseRank, AWP, 보조 GF-PR 점수를 계산한다(감쇠 $d=0.85$, 허용오차 $10^{-8}$, 최대 $300$회 반복).
\item \textbf{프록시, 정렬, 벤치마크:} 제~\ref{ch:methodology}장에서 정의한 다섯 프록시 계열을 계산하고, 각 점수 벡터에 대한 Spearman $\rho$와 Kendall $\tau$를 보고하며, 사전 구축된 간선에 대한 PageRank 해법 시간을 측정하고(3회 평균), 감쇠, 상위 토큰, 표본 크기에 대한 강건성 스윕을 수행한다.
\end{enumerate}

오프라인 합성 검사는 픽스처 그래프에서 동일한 점수·정렬 논리를 사용하여(클라우드 접근 없음) 실제 추출 전에 구현을 검증한다. 실제 추출에는 Google Cloud 자격 증명과 과금 프로젝트가 필요하며, 스캔은 공개 Arbitrum One 로그 데이터셋만을 대상으로 한다.

\dissertationSubheading[sec:bq-extraction]{BigQuery 추출 설계}

ERC-20 \texttt{Approval} 및 \texttt{Transfer} 이벤트는 표준 topic-$0$ 해시로 식별되며, owner/spender 또는 from/to 주소가 평가 지갑 목록에 속하는 로그만으로 제한한다. 따라서 스캔 비용은 전체 체인이 아니라 코호트에 비례한다. 월별 배치는 각 쿼리를 운영상 바이트 상한 아래로 유지한다. GMX V2 이벤트는
\texttt{0xC8ee91A54287DB53897056e12D9819156D3822Fb}
주소의 \texttt{EventEmitter} 계약에서 읽고, \texttt{account}, \texttt{basePnlUsd}, \texttt{sizeDeltaUsd}, 청산(liquidation) 플래그에 대한 키--값 필드로 디코딩한다. 대표 SQL 패턴은 부록~\ref{ch:appendix-sql}에 제시한다.

추출 기록(처리 바이트, 행 수, 타임스탬프)은 감사 가능성을 위해 연구의 수집 로그에 보존한다. 본 논문은 해당 운영 로그가 아니라 그 결과로 얻은 코호트 크기와 간선 수를 보고한다.

\dissertationSubheading[sec:scaling-appendix]{스케일링 표}

표~\ref{tab:benchmark-scaling-incohort}는 $n \leq 5{,}521$에 대한 코호트 내 런타임 스케일링을 보고한다. 제~\ref{ch:results}장의 강건성 표는 주 평가와 동일한 솔버 설정 및 코호트 정의를 사용한다.

% Auto-generated evaluation table fragment — regenerate from the evaluation pipeline; do not edit by hand.

\begin{table}[htbp]
\centering
\caption{In-cohort runtime scaling (SHA256 seed \texttt{benchmark-scale-v1}; $n \leq 5{,}521$ matched GMX cohort).}
\label{tab:benchmark-scaling-incohort}
\fitwidth{%
\begin{tabular}{rrrrrr}
\toprule
$n$ & ER (s) & AWP (s) & Speedup & ER $|E|$ & AWP $|E|$ \\
\midrule
1000 & 0.393 & 3.215 & 8.2$\times$ & 2584 & 23058 \\
2000 & 0.729 & 6.794 & 9.3$\times$ & 5147 & 47031 \\
4000 & 2.075 & 18.031 & 8.7$\times$ & 10557 & 97105 \\
5521 & 2.105 & 18.711 & 8.9$\times$ & 14727 & 137087 \\
\bottomrule
\end{tabular}
}
\end{table}


\dissertationSubheading[sec:delegation-deferred]{신용 위임 추출(보류)}

Aave V3 \texttt{BorrowAllowanceDelegated} 이벤트는 본 연구에서 추출하지 않았다. 따라서 신용 위임 간선에 의존하는 탐색적 순위는 현재 표본에서 정의되지 않으며, 아카이브된 Aave 지향 행렬(부록~\ref{ch:appendix-archive})에서 해석 불가($\tau = \text{---}$)로 보고한다. 위임 기반 평판으로 연구를 확장하려면 해당 이벤트의 전용 추출과 새로운 정렬 평가가 필요하며, 이는 향후 연구로 남긴다.

\dissertationSubheading[sec:supplementary-summary]{보조 정렬 요약 표}

표~\ref{tab:alignment-summary}는 표~\ref{tab:alignment-three-methods}의 계열별 평균 $\tau$를 열 형식(EndorseRank / AWP / GF-PR / 최고 정렬 라벨)으로 다시 제시한다. 본문 서술은 표~\ref{tab:alignment-three-methods}를 사용하며, 열 형식은 나란히 읽기 위해 여기에 유지한다.

% Real BigQuery data -- regenerate with:
%   A-Skill-Programs/margin_rank/scripts/run_dissertation_eval.py --real --export-latex
% Generated by export_latex_results.py -- do not edit by hand unless re-export fails
\begin{table}[htbp]
\centering
\caption[Mean Kendall $\tau$ by proxy family: EndorseRank and AWP.]{Mean Kendall $\tau$ by proxy family: EndorseRank and AWP ($n=5{,}521$).}
\label{tab:alignment-summary}
\fitwidth{%
\begin{tabular}{lccc}
\toprule
Family & EndorseRank & AWP & Higher \\
\midrule
Transfer & 0.333 & 0.534 & AWP \\
Allowance & 0.355 & -0.030 & EndorseRank \\
Liq. & 0.093 & 0.113 & AWP \\
Inv.risk & 0.035 & 0.107 & AWP \\
Sybil & 0.233 & 0.286 & AWP \\
GMX & 0.096 & 0.179 & AWP \\
\bottomrule
\end{tabular}
}
\end{table}


\dissertationSubheading[sec:parameter-summary]{매개변수 요약}

\begin{table}[htbp]
\centering
\caption{주 평가에 사용된 기본 매개변수.}
\label{tab:parameter-summary}
\fitwidth{%
\begin{tabular}{lll}
\toprule
매개변수 & 값 & 역할 \\
\midrule
PageRank 감쇠 $d$ & $0.85$ & 기본 텔레포트 혼합 \\
수렴 허용오차 $\varepsilon$ & $10^{-8}$ & 거듭제곱 반복 중단 기준 \\
최대 반복 횟수 & $300$ & 거듭제곱 반복 상한 \\
로지스틱 $k$ & $0.01$ & AWP 감쇠 기울기 \\
로지스틱 $t_0$ & $180$일 & AWP 감쇠 중점 \\
최소 GMX close 수 & $3$ & 코호트 적격성 \\
관측 구간 & 2025-12-01 -- 2026-05-31 & UTC 양끝 포함 \\
GF-PR 손실 $\varepsilon_{\mathrm{loss}}$ & $0.1$ & 일반 손실 간선 가중 증가 \\
GF-PR liquidation $\varepsilon_{\mathrm{liq}}$ & $1.0$ & 청산(liquidation) 간선 가중 증가 \\
확장 풀 부분표본 시드 & \texttt{benchmark-tier2-v1} & 결정적 SHA256 시드 \\
\bottomrule
\end{tabular}
}
\end{table}

\dissertationSubheading[sec:proxy-appendix]{프록시 운영 정의(참고)}

다음 공식은 제~\ref{ch:methodology}장과 일치한다.

\begin{itemize}
\item \textbf{Transfer:} $\mathrm{in\_degree}(w)=|\{u:\exists u\to w\}|$; $\mathrm{in\_value}(w)=\sum_{e:u\to w}\mathrm{value}(e)$.

\item \textbf{Allowance:} $\mathrm{in\_approve\_degree}(w)=|\{o:\exists o\to w\}|$; $\mathrm{in\_approve\_value}(w)=\sum_o w_a(o,w)$ ($w$는 spender).

\item \textbf{Inverse risk:} $\mathrm{loss\_avoidance}(w)=-\sum_c\min(\mathrm{PnL}_c,0)$; $\mathrm{non\_loss\_close\_rate}(w)=1-|\{c:\mathrm{PnL}_c<0\land\neg\mathrm{liq}_c\}|/|\mathcal{C}(w)|$; $\mathrm{worst\_close\_pnl\_score}(w)=-\min_c\mathrm{PnL}_c$.

\item \textbf{Sybil-adjusted stability:} $\mathrm{inbound\_counterparty\_ratio}(w)=|\{u:u\to w\}|/|T^{\mathrm{in}}(w)|$ (자기 전송 제외); 보유 기간(일); 활성 월 수.

\item \textbf{Trading success:} $\mathrm{close\_success\_count}(w)=|\{c:\mathrm{PnL}_c>0\land\neg\mathrm{liq}_c\}|$; $\mathrm{realized\_gain\_proxy}(w)=\sum_c\max(\mathrm{PnL}_c,0)$; $\mathrm{close\_success\_rate}(w)=\mathrm{close\_success\_count}(w)/|\mathcal{C}(w)|$.
\end{itemize}

\dissertationSubheading[sec:glossary]{용어집}

\begin{description}
\item[Allowance] 소유자를 대신하여 spender가 토큰을 이전할 수 있도록 하는 ERC-20 승인.
\item[Alignment] 평판 점수와 검증 프록시 간의 순위 일치도($\rho$, $\tau$).
\item[Construct overlap] 순위 방법과 프록시가 데이터 출처를 공유하여 정렬을 부풀리는 현상.
\item[Endorsement graph] 최신 allowance의 방향 그래프(EndorseRank 입력).
\item[Matched cohort] 적격성 규칙을 충족하는 주 표본 $n=5{,}521$개 지갑.
\item[Expanded wallet pool] 런타임 스케일링을 위한 더 큰 주소 집합($N=31{,}612$).
\item[PageRank] 감쇠 무작위 보행의 정상분포에 의한 스펙트럼 순위.
\item[Proxy] 영역 직관 검증에 사용되는 관측 가능한 지갑 지표.
\item[Teleport] 확률 $1-d$로 이루어지는 PageRank의 균등 점프.
\end{description}

% !TEX root = ../main.tex
\section{SQL 템플릿, 설정, 확장 주석}
\label{ch:appendix-sql}

본 부록은 본문을 복잡하게 만들지 않으면서 재현을 지원하는 대표 SQL 패턴, 설정 키, 확장된 방법론 주석을 모은다.

\dissertationSubheading[sec:sql-approval]{Approval 로그 추출 패턴}

ERC-20 \texttt{Approval} 이벤트는 topic~$0$이 \texttt{Approval(address,address,uint256)}의 keccak256 해시와 같을 때 식별된다. 단순화된 월별 쿼리 패턴은 다음과 같다.

\begin{verbatim}
SELECT
  block_timestamp,
  block_number,
  transaction_hash,
  log_index,
  address AS token,
  topics[SAFE_OFFSET(1)] AS owner_topic,
  topics[SAFE_OFFSET(2)] AS spender_topic,
  data AS value_data
FROM `bigquery-public-data.goog_blockchain_arbitrum_one_us.logs`
WHERE block_timestamp >= TIMESTAMP('2025-12-01')
  AND block_timestamp <  TIMESTAMP('2026-01-01')
  AND topics[SAFE_OFFSET(0)] = '0x8c5be1e5ebec7d5bd14f71427d1e84f3dd0314c0f7b2291e5b200ac8c7c3b925'
  AND (
    LOWER(CONCAT('0x', SUBSTR(topics[SAFE_OFFSET(1)], 27))) IN UNNEST(@wallet_list)
    OR LOWER(CONCAT('0x', SUBSTR(topics[SAFE_OFFSET(2)], 27))) IN UNNEST(@wallet_list)
  );
\end{verbatim}

\texttt{@wallet\_list}를 통한 지갑 필터링은 스캔을 전체 체인이 아니라 코호트에 비례하도록 유지한다. 월별 배치는 연구의 수집 기록에 문서화된 쿼리당 바이트 상한을 준수한다.

\dissertationSubheading[sec:sql-transfer]{Transfer 로그 추출 패턴}

\texttt{Transfer} 이벤트는 \texttt{Transfer(address,address,uint256)}에 대한 topic~$0$을 사용한다.

\begin{verbatim}
SELECT
  block_timestamp,
  block_number,
  transaction_hash,
  log_index,
  address AS token,
  topics[SAFE_OFFSET(1)] AS from_topic,
  topics[SAFE_OFFSET(2)] AS to_topic,
  data AS value_data
FROM `bigquery-public-data.goog_blockchain_arbitrum_one_us.logs`
WHERE block_timestamp >= TIMESTAMP('2025-12-01')
  AND block_timestamp <  TIMESTAMP('2026-01-01')
  AND topics[SAFE_OFFSET(0)] = '0xddf252ad1be2c89b69c2b068fc378daa952ba7f163c4a11628f55a4df523b3ef'
  AND (
    LOWER(CONCAT('0x', SUBSTR(topics[SAFE_OFFSET(1)], 27))) IN UNNEST(@wallet_list)
    OR LOWER(CONCAT('0x', SUBSTR(topics[SAFE_OFFSET(2)], 27))) IN UNNEST(@wallet_list)
  );
\end{verbatim}

후처리에서는 토큰 소수점을 적용하고, 시간 감쇠 간선 가중치를 구축하며, PageRank 이전에 노드를 평가 표본으로 제한한다.

\dissertationSubheading[sec:sql-gmx]{GMX PositionDecrease 추출 패턴}

GMX V2 이벤트는 \texttt{EventEmitter} 계약에서 읽는다. 파이프라인은 이미터 주소
\texttt{0xC8ee91A54287DB53897056e12D9819156D3822Fb}
로 필터링하고, 키--값 필드를 디코딩하여 \texttt{account}, \texttt{basePnlUsd}, \texttt{sizeDeltaUsd}, 청산(liquidation) 플래그를 복원한다. 프록시 계산에는 \texttt{PositionDecrease} 행만 들어간다. close가 세 번 미만인 지갑은 매칭 코호트에서 제외한다.

\dissertationSubheading[sec:config-keys]{설정 키}

파이프라인 매개변수는 단일 버전 관리 설정 파일에 모인다. 주요 키는 다음과 같다.

\begin{itemize}
\item 관측 구간 시작·종료 타임스탬프;
\item PageRank 감쇠, 허용오차, 최대 반복 횟수;
\item 로지스틱 감쇠 $k$와 $t_0$;
\item GMX 이미터 주소 및 topic 필터;
\item 코호트 적격성을 위한 최소 close 수($m=3$);
\item 벤치마크 반복 횟수 및 부분표본 시드;
\item 중간 산출물과 LaTeX 조각의 출력 위치.
\end{itemize}

감쇠(damping) 또는 로지스틱 감쇠 매개변수를 변경하면 강건성 스윕과 표 내보내기를 포함한 전체 평가를 다시 실행하고, 동일한 산출물에서 그림을 재생성해야 보고된 수치가 파이프라인의 기계 판독 가능 요약과 일치한다.

\dissertationSubheading[sec:complexity-notes]{복잡도 주석}

노드 수를 $n$, 간선 수를 $m=|E|$라 하자. 최신 allowance 간선 구축은 $R_a$개의 approval 행을 타임스탬프로 정렬하면 $O(R_a\log R_a)$이고, (token, owner, spender)를 키로 하는 해시맵을 쓰면 기대 시간 $O(R_a)$이다. 전송 간선 구축은 전송 행에 대해 $O(R_t)$이다. PageRank는 $I$회 반복에 대해 $O(mI)$이다. 경험적으로 매칭 코호트에서 EndorseRank는 $I\approx 37$, AWP는 $I\approx 100$이며 $m_{\mathrm{approve}}\ll m_{\mathrm{transfer}}$이므로, 두 요인 모두 EndorseRank에 유리하다.

메모리는 간선 목록과 점수 벡터가 지배한다. 최대 해법 메모리 $4.6$\,MB 대 $38.2$\,MB(표~\ref{tab:benchmark-runtime})는 간선 수 비 $14{,}727:137{,}087$와 일치한다.

\dissertationSubheading[sec:interpretation-checklist]{독자를 위한 해석 체크리스트}

제~\ref{ch:results}장 표를 읽을 때, 다음 체크리스트는 흔한 오독을 줄인다.

\begin{enumerate}
\item 행이 \emph{사회적} 방법(EndorseRank, AWP)인지 보조 GF-PR 참조인지 식별한다.
\item 각 계열을 의도된 구성체(allowance 대 transfer 대 결과)에 대응시킨다.
\item 절대 $\tau$ 임계값보다 동일 코호트에서의 상대 비교를 우선한다.
\item 거래 성공 및 역위험에서의 GF-PR 고값을 구성체 중첩 진단(SRQ1)으로 취급한다.
\item 런타임 표를 사전 구축된 간선에 대한 PageRank 해법 시간으로 읽고, 종단 간 ETL 시간으로 읽지 않는다.
\item 점 추정치를 일반화하기 전에 강건성 표를 확인한다.
\end{enumerate}

\dissertationSubheading[sec:archived-baselines]{아카이브된 확장 베이스라인}

7-방법 및 6-Aave 정렬 행렬은 평가 파이프라인과 함께 아카이브되며, 스프레드시트 검사용 CSV 행렬도 함께 보관된다. 여러 변형이 추가 이벤트 유형(예: 신용 위임)을 요구하거나 주된 EndorseRank--AWP 질문 밖의 구성체를 도입하기 때문에 본문 서술에서는 제외한다. 본 연구를 확장하는 연구자는 문서화된 프리셋이 있는 전용 내보내기 유틸리티로 내보내기를 다시 활성화할 수 있다.

\dissertationSubheading[sec:limitations-reprise]{한계의 재진술}

부록 자료는 제~\ref{ch:results}장을 넘어 실증 주장을 확장하지 않는다. 특히 SQL 템플릿은 다른 체인을 조회했음을 함의하지 않으며, 아카이브된 베이스라인은 본 논문에서 동료 비교를 구성하지 않는다. 권위 있는 수치는 매칭 코호트($n=5{,}521$)와 확장 지갑 풀($N=31{,}612$)에 대한 파이프라인의 기계 판독 가능 요약 산출물에 남아 있다.

% !TEX root = ../main.tex
\section{PageRank와 순위 상관의 기초 성질}
\label{ch:appendix-theory}

본 부록은 제~\ref{ch:literature}--\ref{ch:results}장에서 암묵적으로 사용된 기초 성질을 기록한다. 내용은 표준적\citep{49,40,54,36}이며, 자족적 독서를 위해 포함한다.

\dissertationSubheading[sec:pr-existence]{PageRank의 존재와 유일성}

$P$를 $n$개 노드 위의 행-확률 행렬이라 하고, 식~\eqref{eq:transition}에서와 같이 dangling-node 보정을 거친 경우도 포함한다. $d\in(0,1)$에 대해 Google 행렬
\[
G_{\mathrm{PR}}=d P^{\top}+\frac{1-d}{n}\mathbf{1}\mathbf{1}^{\top}
\]
은 양행렬이며 열-확률이다(식~\eqref{eq:google-matrix}의 전치 규약까지 고려). 양행렬에 대한 Perron--Frobenius 정리에 의해, $G_{\mathrm{PR}}$는 절댓값이 1인 고유값을 유일하게 가지며, $\sum_i\pi_i=1$로 정규화된 엄밀히 양의 고유벡터 $\boldsymbol{\pi}$를 갖는다. 따라서 감쇠와 dangling-node 규칙이 고정되면 임의의 방향 그래프에 대해 PageRank 점수는 존재하고 유일하다\citep{40}.

이 유일성은 EndorseRank와 AWP 비교의 토대이다. 매개변수가 고정되면 각 방법은 주어진 간선 집합에서 단일 점수 벡터를 유도한다. 따라서 $\mathbf{s}_{\mathrm{ER}}$와 $\mathbf{s}_{\mathrm{AWP}}$의 차이는 솔버의 다중성이 아니라 간선 의미와 가중치에 기인한다.

\dissertationSubheading[sec:pr-continuity]{간선 가중치에 대한 연속성}

고정된 $d\in(0,1)$에 대해 PageRank는 $P$의 성분에 대한 연속 함수이다. 따라서 allowance 또는 transfer 가중치의 작은 섭동은 $\|\cdot\|_1$ 노름에서 점수의 작은 섭동을 유도하며, Lipschitz 상수는 $d$와 $n$에 의존한다\citep{40}. 연속성은 감쇠를 변화시키거나 상위 토큰으로 제한하는 강건성 검사를 정당화한다. 정렬이 칼날처럼 취약하다면, 작은 하이퍼매개변수 변화가 계열별 $\tau$를 뒤섞을 것이다. 경험적으로 EndorseRank의 allowance $\tau$는 $d\in\{0.75,0.85,0.95\}$에 걸쳐 $\approx 0.355$로 유지되며(표~\ref{tab:robustness-damping}), 안정적인 순위 기하와 일치한다.

\dissertationSubheading[sec:pr-sparsity]{희소성과 반복당 비용}

희소 행렬--벡터 곱 $P^{\top}\mathbf{x}$는 비영 간선 $m$개에 대해 $O(m)$ 산술 연산이 든다. 따라서 거듭제곱 반복은 $I$회에 대해 $O(mI)$이다. 매칭 코호트에서 $m_{\mathrm{ER}}=14{,}727$, $m_{\mathrm{AWP}}=137{,}087$이므로, 동일한 $I$에서도 전송 방법은 반복당 대략 $9.3\times$ 더 비싸다. 관측된 반복 횟수($I_{\mathrm{ER}}\approx 37$, $I_{\mathrm{AWP}}\approx 100$)는 격차를 더 벌리며, 측정된 벽시계 시간 비 $\approx 8.9\times$와 일치한다(표~\ref{tab:benchmark-runtime}).

\dissertationSubheading[sec:spearman-properties]{Spearman \texorpdfstring{$\rho$}{rho}의 성질}

Spearman의 $\rho$는 순위에 적용한 Pearson 상관이다. 어느 한쪽 변수의 엄밀히 증가하는 변환에 불변이며 $[-1,1]$에 놓인다. 완벽한 단조 일치는 $\rho=1$, 완벽한 역일치는 $\rho=-1$을 준다. 동점은 중간순위 조정이 필요하며, 평가 파이프라인은 표준적인 동점 인식 구현을 사용한다. 평판 점수와 프록시는 연속적이고 두꺼운 꼬리를 가지므로, 원시 값에 대한 Pearson 상관보다 순위 기반 연관이 더 적절하다\citep{16}.

\dissertationSubheading[sec:kendall-properties]{Kendall \texorpdfstring{$\tau$}{tau}의 성질}

Kendall의 $\tau$는 일치하는 쌍 순서와 불일치하는 쌍 순서의 정규화된 차이와 같다. 확률적 해석이 가능하다. 서로 다른 두 지갑을 무작위로 고를 때, $\tau$는 두 순위 모두에서 상대 순서가 일치할 확률에서 불일치할 확률을 뺀 값이다\citep{36}. 이 해석은 실무자에게 결과를 전달할 때 유용하다. $\tau=0.355$는 allowance 프록시 순서가 EndorseRank 순서와 명확하지만 불완전하게 일치하는 경향을 의미한다.

큰 $n$에서 $\tau$와 $\rho$는 부호와 크기가 종종 비슷하지만 동일하지는 않다. 본 논문은 Do et al.\citep{16}을 따라 둘 다 보고하며, 계열 평균 $\tau$를 주된 요약으로 사용한다.

\dissertationSubheading[sec:construct-validity-note]{구성체 타당도 주석}

Cronbach와 Meehl\citep{14}은 준거 관련 타당도와 구성체 타당도를 구분한다. 금표준 기본값 라벨이 없는 상황에서, 본 논문은 구성체 지향 전략을 취한다. 각 방법은 의도된 의미를 조작화한 프록시(EndorseRank에 대한 allowance 수신; AWP에 대한 인바운드 전송 활동)와, 인접 구성체를 탐색하는 이차 계열(위험, 안정성, 거래 성공)에 대해 검증된다. 구성체 내 높은 $\tau$는 해석을 지지하고, 구성체 간 낮은 $\tau$는 자동으로 실패가 아니다---판별 타당도를 나타낼 수 있다. GF-PR의 극단적인 거래 성공 $\tau$는 준거가 측정과 입력을 공유하는 \emph{오염된} 준거의 경고 사례이다(SRQ1).

\dissertationSubheading[sec:numerical-summary-sheet]{수치 요약표}

빠른 참조를 위해, 주된 매칭 코호트 수치는 다음과 같다.

\begin{center}
\begin{tabular}{ll}
\toprule
양 & 값 \\
\midrule
매칭 지갑 $n$ & $5{,}521$ \\
GMX close(디코딩) & $365{,}488$ \\
EndorseRank 간선 & $14{,}727$ \\
AWP 간선 & $137{,}087$ \\
EndorseRank 런타임 & $2.105$\,s \\
AWP 런타임 & $18.711$\,s \\
런타임 비(AWP/ER) & $\approx 8.9\times$ \\
방법 간 $\tau$ & $0.316$ \\
ER allowance 평균 $\tau$ & $0.355$ \\
AWP transfer 평균 $\tau$ & $0.534$ \\
ER 거래 성공 평균 $\tau$ & $0.096$ \\
AWP 거래 성공 평균 $\tau$ & $0.179$ \\
GF-PR 거래 성공 평균 $\tau$ & $0.583$ \\
GF-PR 역위험 평균 $\tau$ & $0.376$ \\
확장 풀 $N$ & $31{,}612$ \\
$n=10{,}000$에서의 가속 & $\approx 10.4\times$ \\
$N=31{,}612$에서의 가속 & $\approx 9.4\times$ \\
관측 구간 & 2025-12-01 -- 2026-05-31 \\
\bottomrule
\end{tabular}
\end{center}

이 값들은 파이프라인의 기계 판독 가능 평가 요약에서 전사되었으며, 내보낸 결과 표와 일치한다.

% !TEX root = ../main.tex
\section{아카이브된 확장 베이스라인}
\label{ch:appendix-archive}

탐색적 분석 과정에서, 동일한 매칭 코호트와 프록시 계열에 대해 추가 PageRank 변형을 점수화하였다. 완전성과 향후 연구를 위해 여기에 아카이브한다. 이들은 주된 EndorseRank--AWP 연구 질문의 일부가 \emph{아니며}, 본문 서술에서 동료 주장으로 읽혀서는 안 된다.

\dissertationSubheading[sec:seven-method-archive]{7-방법 정렬 행렬}

표~\ref{tab:alignment-seven-methods}는 LP-PR, CW-AWP, LF-PR, RiskProp 스타일 변형을 포함하여 평가 코드베이스에 구현된 일곱 방법에 대한 계열별 평균 Kendall $\tau$를 보고한다. EndorseRank는 allowance 계열의 선두를 유지하고, AWP(및 밀접한 전송 변형)는 transfer에서 선두이며, GF-PR는 구성상 역위험과 GMX 성공에서 선두이다.

% Real BigQuery data -- regenerate with:
%   run_dissertation_eval.py --real --export-latex

\begin{table}[htbp]
\centering
\caption{Mean Kendall $\tau$ by proxy family across 7 PageRank methods ($n=5521$). Bold = column maximum. Runtime = mean PageRank wall time (s, 5 timed repeats after one warm-up). Edge counts: liq_pr=45, borrow_pr=484, repay_pr=787, borrow_pr_pool=262, repay_pr_pool=232, delegation_pr=0.}
\label{tab:alignment-seven-methods}
\fitwidth{%
\begin{tabular}{lccccccr}
\toprule
Method & Transfer & Allowance & Liq. & Inv.risk & Sybil & GMX & Runtime (s) \\
\midrule
AWP & 0.534 & -0.030 & 0.113 & 0.107 & \textbf{0.286} & 0.179 & 11.253 \\
EndorseRank & 0.333 & \textbf{0.355} & 0.093 & 0.035 & 0.233 & 0.096 & 1.426 \\
GF-PR & 0.246 & -0.041 & 0.132 & \textbf{0.376} & 0.093 & \textbf{0.583} & 0.695 \\
LP-PR & 0.529 & -0.027 & \textbf{0.136} & 0.098 & 0.280 & 0.196 & 10.377 \\
CW-AWP & 0.504 & -0.013 & 0.128 & 0.114 & 0.253 & 0.204 & 9.776 \\
LF-PR & 0.197 & -0.031 & 0.045 & 0.017 & 0.102 & 0.032 & 0.064 \\
RiskProp & \textbf{0.534} & -0.030 & 0.113 & 0.107 & 0.286 & 0.179 & 13.453 \\
\bottomrule
\end{tabular}
}
\end{table}


해석 지침:

\begin{itemize}
\item transfer형 간선에 구축된 방법은 transfer 및 Sybil 계열에서 AWP 근처에 군집한다.
\item 간선 수가 매우 작은 방법(예: 수십 개 간선의 liquidation 지향 그래프)은 불안정하거나 퇴화한 순위를 낼 수 있으며, 런타임이 낮은 것은 단순히 $m$이 작기 때문일 수 있다.
\item AWP를 수치적으로 반영하는 RiskProp 스타일 행은 독립적 발견이 아니라 공유된 전송 구조를 나타낸다.
\item 결과 계열에서 GF-PR의 우세는 제~\ref{ch:results}장과 동일한 구성체 중첩 경고로 남는다.
\end{itemize}

\dissertationSubheading[sec:six-aave-archive]{6-방법 Aave 지향 행렬}

표~\ref{tab:alignment-six-aave-methods}는 Aave 지향 프리셋을 보고한다. 여러 대출 이벤트 그래프는 연구 구간 내 Arbitrum에서 희소하며, \texttt{BorrowAllowanceDelegated} 이벤트가 추출되지 않은 경우 위임 기반 행은 퇴화할 수 있다(부록~\ref{ch:appendix-eval}).

% Real BigQuery data -- regenerate with:
%   run_dissertation_eval.py --real --export-latex

\begin{table}[htbp]
\centering
\caption{Mean Kendall $\tau$ by proxy family across 6 PageRank methods ($n=5521$). Bold = column maximum. Runtime = mean PageRank wall time (s, 5 timed repeats after one warm-up). Aave W\textrightarrow W edges: LiqCall (liquidator$\rightarrow$user), Borrow (initiator$\rightarrow$onBehalfOf), Repay (repayer$\rightarrow$user), Delegation (delegator$\rightarrow$delegatee). Credit Delegation BQ extract deferred when parquet absent---$\tau$ may be degenerate. Edge counts: \texttt{liq_pr=45}, \texttt{borrow_pr=484}, \texttt{repay_pr=787}, \texttt{borrow_pr_pool=262}, \texttt{repay_pr_pool=232}, \texttt{delegation_pr=0}.}
\label{tab:alignment-six-aave-methods}
\fitwidth{%
\begin{tabular}{lccccccr}
\toprule
Method & Transfer & Allowance & Liq. & Inv.risk & Sybil & GMX & Runtime (s) \\
\midrule
AWP & \textbf{0.534} & -0.030 & \textbf{0.113} & \textbf{0.107} & \textbf{0.286} & \textbf{0.179} & 11.253 \\
EndorseRank & 0.333 & \textbf{0.355} & 0.093 & 0.035 & 0.233 & 0.096 & 1.426 \\
LiqCall-PR & 0.053 & -0.012 & -0.008 & 0.026 & 0.025 & 0.030 & 0.009 \\
Borrow-PR & 0.113 & -0.019 & 0.020 & 0.014 & 0.055 & 0.023 & 0.056 \\
Repay-PR & 0.143 & -0.024 & 0.016 & 0.022 & 0.068 & 0.026 & 0.082 \\
Delegation-PR & --- & --- & --- & --- & --- & --- & 0.000 \\
\bottomrule
\end{tabular}
}
\end{table}


이 행렬들은 신용 위임 간선과 교차 프로토콜 대출 라벨에 대한 향후 연구를 동기화하지만, 주장 C1--C4 또는 SC1을 변경하지는 않는다.

\dissertationSubheading[sec:why-archived]{본문이 세 방법을 사용하는 이유}

논문을 EndorseRank, AWP, 보조 GF-PR에 집중시키는 것은 세 가지 목표를 위한 것이다.

\begin{enumerate}
\item \textbf{기여의 명확성:} 새로운 구성체는 allowance 기반 보증이다. 확립된 전송 베이스라인\citep{16}과 비교하면 연구 질문이 분리된다.
\item \textbf{베이스라인 난립 회피:} AWP와 전송 간선을 공유하는 추가 변형은 독립 구성체를 더하지 않고 행만 늘린다.
\item \textbf{GF-PR의 정직한 사용:} GF-PR를 보조로 두면, 독자가 결과 고유의 상한을 사회적 평판 경쟁자로 취급하는 것을 막는다.
\end{enumerate}

더 넓은 베이스라인이 필요한 연구자는 공개 평가 파이프라인에서 7-방법 및 6-Aave 프리셋을 사용하여 아카이브 표를 재생성할 수 있다.

\dissertationSubheading[sec:archive-repro]{아카이브 표 재현}

아카이브 표는 평가 파이프라인의 표 내보내기에서 7-방법 및 6-Aave 프리셋으로 재생성된다. CSV 사본은 스프레드시트 검사용으로 LaTeX 조각과 함께 기록된다.

\dissertationSubheading[sec:archive-method-notes]{아카이브 방법 식별자에 대한 주석}

다음 식별자는 아카이브 표와 CSV 내보내기에 나타난다. 평가 코드베이스의 탐색적 구현을 문서화한다.

\begin{description}
\item[LP-PR] 개발 중 탐색한, 대안적 손실 페널티 또는 유동성 지향 가중치를 가진 transfer 그래프 변형. 경험적으로 transfer 계열에서 AWP를 밀접히 추적하며(표~\ref{tab:alignment-seven-methods}), 공유된 결제 흐름 구조를 나타낸다.
\item[CW-AWP] 동일 쌍 간 반복 전송의 가중치를 낮추는 상대방 가중 AWP 변형. AWP 대비 transfer $\tau$를 약간 낮추면서도 동일 구성체 계열에 남는다.
\item[LF-PR] 간선 수가 매우 작고 런타임이 거의 0인 저빈도 또는 약하게 필터링된 그래프. 순위는 전체 AWP와 비교할 수 없으며, 실패한 희소 베이스라인을 문서화하기 위해서만 유지한다.
\item[RiskProp] Lin et al.\citep{44}에서 영감을 받은, transfer형 간선 위의 위험 전파 스타일 순위. 본 코호트에서는 여러 계열에서 AWP를 수치적으로 반영하며, 구현된 위험 신호가 연구의 매개변수화 아래에서 전송 중심성과 강하게 갈라지지 않았음을 시사한다.
\item[LiqCall-PR / Borrow-PR / Repay-PR] 해당 이벤트가 있을 때 liquidation call, borrow, repay로 구축한 Aave 지향 지갑-대-지갑 그래프. 연구 구간 Arbitrum에서 간선 수는 작아(수십에서 수백) 정렬이 약하고 런타임은 무시할 만하다.
\item[Delegation-PR] 신용 위임 allowance를 사용하도록 의도되었으나, 추출된 위임 간선이 0이므로 점수는 퇴화하고 $\tau$는 해석할 수 없다.
\end{description}

이 주석은 향후 연구자가 아카이브 행을 동료 심사된 베이스라인으로 취급하기보다, 어떤 변형이 완전한 재구현을 받을 가치가 있는지 판단할 수 있도록 존재한다.

\dissertationSubheading[sec:archive-decision-log]{결정 로그: 본문에 남은 것}

제~\ref{ch:results}--\ref{ch:discussion}장에 EndorseRank, AWP, 보조 GF-PR만 남기기로 한 결정은 7-방법 및 6-Aave 행렬을 검토한 뒤에 이루어졌다. 기준은 다음과 같다.

\begin{enumerate}
\item \textbf{구성체 독립성:} 베이스라인은 전송의 사소한 재가중이 아니라 구별되는 간선 의미를 구현해야 한다.
\item \textbf{데이터 충분성:} 매칭 코호트에서 안정적인 PageRank를 위해 간선 수가 충분히 커야 한다.
\item \textbf{서술 초점:} 논문의 기여는 allowance 기반 보증이며, 가능한 모든 PageRank 변형의 조사가 아니다.
\item \textbf{평가의 정직성:} 결과 고유 방법은 보조 진단으로 남긴다.
\end{enumerate}

이 기준 아래에서 EndorseRank(allowance), AWP(transfer), GF-PR(결과 고유 상한)는 최소 충분 집합을 이룬다. 아카이브 표는 주된 주장을 희석하지 않으면서 더 넓은 지형을 원하는 독자를 위해 이용 가능하다.

% !TEX root = ../main.tex
\section{주장--증거 대응표}
\label{ch:appendix-claims}

본 부록은 각 가설, 연구 질문, 실증 주장을 제~\ref{ch:results}장의 구체적 증거에 대응시킨다. 심사자와 재현자를 위한 탐색 보조로 의도되었다.

\dissertationSubheading[sec:map-h1]{H1 / RQ1 --- 순위 분기}

\begin{itemize}
\item \textbf{진술:} EndorseRank와 AWP는 동일 코호트에서 구별되는 지갑 순서를 산출한다.

\item \textbf{증거:} 점수 벡터 간 방법 간 Kendall $\tau=0.316$(제~\ref{sec:rank-divergence}절); 그림~\ref{fig:rank-scatter}의 순위 산점도; 그림~\ref{fig:score-distributions}의 점수 분포.

\item \textbf{해석:} 관련되나 상호 교환 가능하지 않은 순위.
\end{itemize}

\dissertationSubheading[sec:map-h2]{H2 / RQ2 --- 구성체별 정렬}

\begin{itemize}
\item \textbf{진술:} 사회적 방법은 정렬하는 프록시 계열에서 서로 다르다.

\item \textbf{증거:} 표~\ref{tab:alignment-three-methods}; 계열별 표~\ref{tab:alignment-transfer}--\ref{tab:alignment-gmx}; 히트맵 그림~\ref{fig:alignment-heatmap}.

\item \textbf{승자:} allowance에서 EndorseRank($\tau=0.355$ 대 $-0.030$); transfer에서 AWP($0.534$ 대 $0.333$), 사회적 방법 중 거래 성공에서 AWP($0.179$ 대 $0.096$).
\end{itemize}

\dissertationSubheading[sec:map-h3]{H3 / RQ3 --- 계산 효율}

\begin{itemize}
\item \textbf{진술:} 동일 $n$에서 EndorseRank가 AWP보다 빠르다.

\item \textbf{증거:} 표~\ref{tab:benchmark-runtime}($2.105$\,s 대 $18.711$\,s, $\approx 8.9\times$); 표~\ref{tab:benchmark-scaling}($n=10{,}000$에서 $\approx 10.4\times$; $N=31{,}612$에서 $\approx 9.4\times$); 그림~\ref{fig:runtime-scaling}; 간선 수 $14{,}727$ 대 $137{,}087$.
\end{itemize}

\dissertationSubheading[sec:map-rq4]{RQ4 --- 검증 틀}

\begin{itemize}
\item \textbf{진술:} 사회적 방법에 대해, 본 체인에서 transfer 프록시가 거래 성공 프록시보다 더 강하게 정렬한다.

\item \textbf{증거:} 표~\ref{tab:alignment-three-methods}: transfer 평균 $0.333$(ER) 및 $0.534$(AWP)가 거래 성공 평균 $0.096$(ER) 및 $0.179$(AWP)를 초과한다.
\end{itemize}

\dissertationSubheading[sec:map-srq1]{SRQ1 / SC1 --- GF-PR 상한}

\begin{itemize}
\item \textbf{진술:} GF-PR는 구성체 중첩을 통해 거래 결과 정렬의 상한을 제공한다.

\item \textbf{증거:} 표~\ref{tab:alignment-three-methods}의 GF-PR 행: 역위험 $0.376$, 거래 성공 $0.583$; 제~\ref{sec:gf-pr-ceiling}절.
\end{itemize}

\dissertationSubheading[sec:map-c1c4]{주장 C1--C4}

\begin{itemize}
\item \textbf{C1:} H3/RQ3와 동일한 런타임 증거.
\item \textbf{C2:} 선택된 프록시에서의 신뢰할 만한 정렬---ER에 대한 allowance, AWP에 대한 transfer---표~\ref{tab:alignment-allowance} 및~\ref{tab:alignment-transfer}.
\item \textbf{C3:} H2/RQ2와 같은 구성체별 승자.
\item \textbf{C4:} C1과 EndorseRank의 allowance $\tau=0.355$ 및 transfer $\tau=0.333$을 결합한 효율--정렬 트레이드오프; 표~\ref{tab:robustness-damping}--\ref{tab:robustness-sample-size} 및 그림~\ref{fig:damping-sensitivity}의 강건성.
\end{itemize}

\dissertationSubheading[sec:map-nonclaims]{명시적 비주장}

본 논문은 다음을 \emph{주장하지 않는다}.

\begin{itemize}
\item EndorseRank 또는 AWP가 사회적 그래프만으로 거래 이익을 예측한다는 것;
\item AWP가 구식이 되었다는 것;
\item GF-PR가 배포 가능한 사회적 평판 점수라는 것;
\item 결과가 재현 없이 Arbitrum One, GMX V2, 연구 구간을 넘어 일반화된다는 것;
\item 온체인 점수가 담보 또는 법적 신원 체계를 대체한다는 것.
\end{itemize}

\dissertationSubheading[sec:map-artifacts]{재현을 위한 산출물 체크리스트}

\begin{enumerate}
\item 병합된 지갑 순위 표 --- ER, AWP, GF-PR 점수.
\item PnL 및 청산(liquidation) 플래그가 있는 디코딩된 GMX close 이벤트.
\item 기계 판독 가능 평가 요약 --- 권위 있는 수치.
\item 내보낸 LaTeX 표 조각.
\item 내보낸 벡터 그림.
\item 선택적 아카이브 확장 베이스라인 행렬.
\end{enumerate}

재현자는 강건성 스윕과 표 내보내기를 포함하여 실제 데이터에서 평가 파이프라인을 실행한 뒤, 동일한 산출물에서 그림을 재생성하여 표와 그림을 다시 만들 수 있다.

